\SourceLeaf{367}%
and a short extract from it was printed in our journal at the time (1805, no.~124). Here we note only that this is the passage in which Socrates uses an example from geometry to show how one should prepare to solve a problem by first adopting certain more precise specifications. The geometrical problem Socrates chooses is whether it is possible to inscribe a given triangle in a given circle. The words by which he first seeks to impose certain restrictions on the kind of triangle, however, have given his interpreters much trouble. Mr. \textsc{Mollweide} shows with much scholarly acumen that the property referred to is the ability to divide the triangle into two other triangles similar to the whole. This is, of course, essentially just an affected way of saying \textit{right-angled} triangle. Mr. M.'s proof that the property can belong only to a right-angled triangle is much more intricate and lengthy than was needed here, since this follows immediately from the equality of the three angles \(ABC,ADB,BDC\) (p.~46).

The third treatise, too, had already been submitted to our Society in manuscript, and an account of it appeared in our \textit{Gelehrte Anzeigen} (1807, no.~74). It appears here with substantial additions. It explains two methods given by \textsc{Columella} for finding approximate areas, one for an equilateral triangle and one for a circular segment. There is a small oversight on p.~71, where it is claimed that no other fraction with both numerator and denominator under 100 can come as close as \(\frac{13}{30}\) to the true ratio of the equilateral triangle to the square on the same side. The two fractions \(\frac{29}{67}\) and \(\frac{42}{97}\) are in fact more accurate.

The letter to the distinguished Professor \textsc{Schneider} in Breslau contains a few remarks on the excerpts that \textsc{Hase} presented in \textsc{Bredow}’s \textit{Epistolae Parisienses} from the admittedly very unimportant mathematical writings of \textsc{Vitruvius Rufus} and \textsc{Epaphroditus}.
